\section{Discussão Técnica e Limitações}

% TODO: Discutir criticamente os resultados observados:
% 1. Compromisso quantitativo entre APF (trajetorias mais suaves e rapidas em obstaculos esparsos) e Bug2 (garantia de convergencia em armadilhas concavas e maior margem d_min, ao custo de maior tempo/comprimento em contornos).
% 2. Interpretacao fisica/matematica da estagnacao do APF no Cenario B (anulacao F_att + F_rep = 0 em x ~ 2,85 m) e no Subcenario GNRON (relacao entre rho_0 e a distancia da parede atras da meta).
% 3. Influencia da odometria RK2 e do ruido sensorial (encoders e LiDAR) sobre a precisao final (e_p(t_f) < 1 cm nas execucoes sem colisao; degradacao acentuada em caso de colisao na configuracao Fraco).
% 4. Limitacoes da abordagem reativa classica (dependencia de estimativa odométrica sem correcao exteroceptiva em percursos longos, comportamento na linha-m sob deriva).
